\documentclass[11pt]{article}
\usepackage[margin=1in]{geometry}
\usepackage{amsmath,amssymb,amsthm}
\usepackage{graphicx}
\usepackage{hyperref}

\newtheorem{theorem}{Theorem}
\newtheorem{proposition}{Proposition}
\newtheorem{lemma}{Lemma}
\newtheorem{corollary}{Corollary}

\newcommand{\E}{\mathbb E}
\newcommand{\Pp}{\mathbb P}
\newcommand{\R}{\mathbb R}
\newcommand{\norm}[1]{\left\lVert #1\right\rVert}
\newcommand{\simequiv}{\simeq}

\title{Hilbert-Valued Characteristic Domination for All Finite \(p\)}
\author{Agent \texttt{agent\_lane\_19}}
\date{2026-06-16}

\begin{document}
\maketitle

\section*{Status}

This packet gives a candidate partial result, likely valid, for Remark 12.10 of
Yaroslavtsev's paper \emph{Local characteristics and tangency of vector-valued
martingales} \cite{Yaroslavtsev2019}.  It solves the Hilbert-valued subcase of
the general-local-martingale characteristic-domination problem for every
\(1\leq p<\infty\).  It is partial relative to the broad source question
because it does not cover non-Hilbert UMD spaces.

The new part of this packet is the range \(1\leq p<2\).  The already packaged
Hilbert \(p=2\) and \(p\geq2\) packets handled the upper range by the Hilbert
square identity and Rosenthal/BDG estimates.  Below \(2\), the missing
substitute is a Laplace-transform comparison for random square functions.

\section*{Source Problem}

The source is arXiv:1907.11588, PDF page 105, Remark 12.10.  The source crop
below states that the general-local-martingale case is open and reduces the
missing accessible-jump issue to a discrete independent-increment comparison.

\begin{center}
\includegraphics[width=\textwidth]{figures/open_problem_crop.png}
\end{center}

The next page adds that, by symmetrization, one may assume the variables are
symmetric, and that even this symmetric discrete case was not known to the
author.

\begin{center}

\includegraphics[width=\textwidth]{figures/open_problem_continuation_crop.png}
\end{center}

In discrete form, the source asks whether independent mean-zero \(X\)-valued
variables \((\xi_n)\), \((\eta_n)\) satisfying
\[
   \sum_n \Pp(\xi_n\in B)
   \leq
   \sum_n \Pp(\eta_n\in B)
   \qquad
   (B\subset X\setminus\{0\}\text{ Borel})
\]
must satisfy
\[
   \E\norm{\sum_n \xi_n}^p
   \leq
   C_{p,X}\E\norm{\sum_n \eta_n}^p .
\]

\section*{Main Result}

\begin{theorem}\label{thm:hilbert-all-p}
Let \(H\) be a separable Hilbert space and \(1\leq p<\infty\).  Let \(M\) and
\(N\) be \(H\)-valued local martingales such that \(N\) is characteristically
dominated by \(M\) in the sense of \cite{Yaroslavtsev2019}.  Then
\[
  \E\sup_{t\geq0}\norm{N_t}^p
  \leq
  C_p\,\E\sup_{t\geq0}\norm{M_t}^p ,
\]
whenever the right hand side is finite.
\end{theorem}

\section*{Idea of the Proof}

The hard range is \(1\leq p<2\).  A naive Rosenthal estimate is poorly matched
there: the term \(\sum_n\E\norm{Z_n}^p\) is too large and is not comparable to
the \(p\)-moment of the sum.  The right object is instead the random square
function
\[
   S_Z=\Bigl(\sum_n \norm{Z_n}^2\Bigr)^{1/2}.
\]
For independent mean-zero Hilbert-valued variables,
\[
   \E\norm{\sum_n Z_n}^p \simequiv_p \E S_Z^p,
   \qquad 1\leq p<\infty.
\]
Thus the discrete problem below \(2\) becomes a scalar question: if the
aggregate law of the nonnegative variables \(Y_n=\norm{\xi_n}^2\) is dominated
by that of \(V_n=\norm{\eta_n}^2\), does
\[
   \E\Bigl(\sum_nY_n\Bigr)^\alpha
   \lesssim_\alpha
   \E\Bigl(\sum_nV_n\Bigr)^\alpha,
   \qquad \alpha=p/2<1,
\]
hold?  Yes.  For \(S=\sum_nY_n\),
\[
   \E S^\alpha
   =
   c_\alpha\int_0^\infty (1-\E e^{-tS})t^{-1-\alpha}\,dt,
\]
and independence gives
\[
   1-\E e^{-tS}=1-\prod_n(1-u_n(t)),
   \qquad
   u_n(t)=\E(1-e^{-tY_n}).
\]
Writing \(U(t)=\sum_nu_n(t)\), one has
\[
   (1-e^{-1})\min\{1,U(t)\}
   \leq
   1-\prod_n(1-u_n(t))
   \leq
   \min\{1,U(t)\}.
\]
The quantity \(U(t)\) depends only on the aggregate law.  Hence aggregate
domination controls every fractional square-function moment with
\(\alpha<1\).  This is the below-\(2\) good-lambda substitute.

\section*{Discrete Hilbert Comparison}

\begin{lemma}[Independent Hilbert sums and square functions]
\label{lem:hilbert-mz}
Let \(H\) be a Hilbert space, \(1\leq p<\infty\), and let
\(Z_1,\ldots,Z_N\) be independent mean-zero \(H\)-valued random variables.
Then
\[
   \E\norm{\sum_{n=1}^N Z_n}^p
   \simequiv_p
   \E\Bigl(\sum_{n=1}^N\norm{Z_n}^2\Bigr)^{p/2}.
\]
Moreover,
\[
   \E\max_{1\leq m\leq N}\norm{\sum_{n=1}^m Z_n}^p
   \simequiv_p
   \E\Bigl(\sum_{n=1}^N\norm{Z_n}^2\Bigr)^{p/2}.
\]
\end{lemma}

\begin{proof}
The maximal estimate is the classical Hilbert-valued discrete
Burkholder-Davis-Gundy inequality, recalled for instance in
\cite{Yaroslavtsev2018}.  Since \(H\) is Hilbert, the Gaussian/radonifying
term in the discrete BDG theorem is exactly
\((\sum_n\norm{Z_n}^2)^{1/2}\).

For the terminal estimate, let \(Z'_1,\ldots,Z'_N\) be an independent copy and
write \(S=\sum_n Z_n\), \(S'=\sum_n Z'_n\), and \(W_n=Z_n-Z'_n\).  By Jensen,
\[
   \E\norm{S}^p
   =
   \E\norm{\E(S-S'\mid Z_1,\ldots,Z_N)}^p
   \leq
   \E\norm{S-S'}^p .
\]
The variables \(W_n\) are independent and symmetric, so their joint law is
invariant under independent sign changes.  Hence, for an auxiliary Rademacher
sequence \((\varepsilon_n)\), the Hilbert-space Khintchine inequality gives
\[
   \E\norm{S-S'}^p
   =
   \E\E_\varepsilon\norm{\sum_n\varepsilon_nW_n}^p
   \simequiv_p
   \E\Bigl(\sum_n\norm{W_n}^2\Bigr)^{p/2}
   \lesssim_p
   \E\Bigl(\sum_n\norm{Z_n}^2\Bigr)^{p/2}.
\]
For the final inequality we used the triangle inequality in
\(L^p(\Omega;\ell_2(H))\).

Conversely,
\[
   \E\norm{S-S'}^p
   \leq
   2^p\E\norm{S}^p,
\]
again by the triangle inequality and the equality in law of \(S\) and \(S'\).
The lower Khintchine inequality and Jensen applied in the Hilbert space
\(\ell_2^N(H)\) give
\[
   \E\norm{S}^p
   \gtrsim_p
   \E\Bigl(\sum_n\norm{W_n}^2\Bigr)^{p/2}
   \geq
   \E\Bigl(\sum_n\norm{\E(W_n\mid Z_1,\ldots,Z_N)}^2\Bigr)^{p/2}.
\]
Since \(\E(W_n\mid Z_1,\ldots,Z_N)=Z_n\), this is the desired lower bound.
\end{proof}

\begin{lemma}[Laplace comparison for fractional square functions]
\label{lem:laplace-square}
Let \(0<\alpha<1\).  For \(i=1,\ldots,N\), let \(Y_i\) and \(V_i\) be
independent nonnegative random variables.  Let \(\mu_i\) and \(\nu_i\) be the
restrictions of their laws to \((0,\infty)\).  If
\[
   \sum_{i=1}^N\mu_i \leq \sum_{i=1}^N\nu_i
   \qquad\text{as measures on }(0,\infty),
\]
then
\[
   \E\Bigl(\sum_{i=1}^N Y_i\Bigr)^\alpha
   \leq
   \frac{1}{1-e^{-1}}\,
   \E\Bigl(\sum_{i=1}^N V_i\Bigr)^\alpha .
\]
\end{lemma}

\begin{proof}
For \(S_Y=\sum_iY_i\), set
\[
   u_i(t)=\E(1-e^{-tY_i})
          =\int_{(0,\infty)}(1-e^{-ty})\,d\mu_i(y),
   \qquad
   U_Y(t)=\sum_i u_i(t).
\]
Then \(0\leq u_i(t)\leq1\), and independence gives
\[
   \E e^{-tS_Y}=\prod_i(1-u_i(t)).
\]
Therefore
\[
   1-\E e^{-tS_Y}
   =
   1-\prod_i(1-u_i(t))
   \leq
   \min\{1,U_Y(t)\},
\]
while
\[
   1-\E e^{-tS_Y}
   \geq
   1-e^{-U_Y(t)}
   \geq
   (1-e^{-1})\min\{1,U_Y(t)\}.
\]
The same statements hold for \(S_V=\sum_iV_i\).  The measure domination implies
\[
   U_Y(t)
   =
   \int_{(0,\infty)}(1-e^{-ty})\,d\sum_i\mu_i(y)
   \leq
   \int_{(0,\infty)}(1-e^{-ty})\,d\sum_i\nu_i(y)
   =
   U_V(t).
\]
Hence, for every \(t>0\),
\[
   1-\E e^{-tS_Y}
   \leq
   \frac{1}{1-e^{-1}}\,(1-\E e^{-tS_V}).
\]
Finally use the standard fractional-moment identity
\[
   \E S^\alpha
   =
   \frac{\alpha}{\Gamma(1-\alpha)}
   \int_0^\infty (1-\E e^{-tS})t^{-1-\alpha}\,dt
   \qquad(0<\alpha<1).
\]
\end{proof}

\begin{proposition}[Discrete aggregate domination in Hilbert spaces below \(2\)]
\label{prop:discrete-hilbert-low}
Let \(H\) be a Hilbert space and \(1\leq p<2\).  Let
\((\xi_n)_{n=1}^N\) and \((\eta_n)_{n=1}^N\) be independent mean-zero
\(H\)-valued random variables.  If
\[
   \sum_{n=1}^N \Pp(\xi_n\in B)
   \leq
   \sum_{n=1}^N \Pp(\eta_n\in B)
   \qquad
   (B\subset H\setminus\{0\}\text{ Borel}),
\]
then
\[
   \E\norm{\sum_{n=1}^N \xi_n}^p
   \leq
   C_p\,\E\norm{\sum_{n=1}^N \eta_n}^p .
\]
The same comparison holds with terminal sums replaced by the corresponding
maxima over partial sums in any fixed order.
\end{proposition}

\begin{proof}
Put \(Y_n=\norm{\xi_n}^2\), \(V_n=\norm{\eta_n}^2\), and
\(\alpha=p/2\in[1/2,1)\).  Applying the map \(x\mapsto\norm{x}^2\) to the
aggregate jump measures gives
\[
   \sum_n \mathcal L(Y_n)|_{(0,\infty)}
   \leq
   \sum_n \mathcal L(V_n)|_{(0,\infty)} .
\]
Lemma \ref{lem:laplace-square} yields
\[
   \E\Bigl(\sum_n\norm{\xi_n}^2\Bigr)^{p/2}
   \leq
   C\,\E\Bigl(\sum_n\norm{\eta_n}^2\Bigr)^{p/2}.
\]
The terminal and maximal assertions now follow from Lemma
\ref{lem:hilbert-mz}.
\end{proof}

\begin{proposition}[Discrete aggregate domination in Hilbert spaces for all \(p\)]
\label{prop:discrete-hilbert-all}
Let \(H\) be a Hilbert space and \(1\leq p<\infty\).  Let
\((\xi_n)_{n=1}^N\) and \((\eta_n)_{n=1}^N\) be independent mean-zero
\(H\)-valued random variables satisfying the aggregate domination hypothesis
of Proposition \ref{prop:discrete-hilbert-low}.  Then
\[
   \E\norm{\sum_{n=1}^N \xi_n}^p
   \leq
   C_p\,\E\norm{\sum_{n=1}^N \eta_n}^p ,
\]
and the same comparison holds for maxima over partial sums in any fixed order.
\end{proposition}

\begin{proof}
The case \(1\leq p<2\) is Proposition \ref{prop:discrete-hilbert-low}.  Let
\(p\geq2\) and set \(r=p/2\geq1\).  Lemma \ref{lem:hilbert-mz} reduces matters
to square functions.  For \(Y_n=\norm{\xi_n}^2\), the scalar Rosenthal
inequality for independent nonnegative variables gives
\[
  \E\Bigl(\sum_nY_n\Bigr)^r
  \lesssim_p
  \Bigl(\sum_n\E Y_n\Bigr)^r+\sum_n\E Y_n^r .
\]
Aggregate domination controls the two terms:
\[
   \sum_n\E\norm{\xi_n}^2\leq \sum_n\E\norm{\eta_n}^2,
   \qquad
   \sum_n\E\norm{\xi_n}^p\leq \sum_n\E\norm{\eta_n}^p.
\]
For \(V_n=\norm{\eta_n}^2\),
\[
  \Bigl(\sum_n\E V_n\Bigr)^r+\sum_n\E V_n^r
  \leq
  2\,\E\Bigl(\sum_nV_n\Bigr)^r,
\]
by Jensen and the pointwise inequality \((\sum_nV_n)^r\geq\sum_nV_n^r\).
Thus the \(\xi\)-square-function \(p\)-moment is controlled by the
\(\eta\)-square-function \(p\)-moment, and Lemma \ref{lem:hilbert-mz} finishes
the proof.  The endpoint \(p=2\) is also immediate from the Hilbert square
identity.
\end{proof}

\section*{Independent-Increment Characteristic Comparison}

\begin{proposition}[Hilbert independent-increment comparison]
\label{prop:ind-incr-hilbert}
Let \(H\) be a separable Hilbert space and \(1\leq p<\infty\).  Let \(L\) and
\(K\) be \(H\)-valued martingales with independent increments and deterministic
local characteristics \((A_L,\nu_L)\) and \((A_K,\nu_K)\).  Assume
\[
   \operatorname{tr} A_L(\infty)\leq \operatorname{tr} A_K(\infty),
   \qquad
   \nu_L(\mathbb R_+\times B)\leq \nu_K(\mathbb R_+\times B)
\]
for every Borel \(B\subset H\setminus\{0\}\).  Then
\[
   \E\sup_{t\geq0}\norm{L_t}^p
   \leq
   C_p\,\E\sup_{t\geq0}\norm{K_t}^p .
\]
\end{proposition}

\begin{proof}
By the Hilbert-valued BDG inequality,
\[
   \E\sup_{t\geq0}\norm{L_t}^p
   \simequiv_p
   \E Q_L^{p/2},
   \qquad
   Q_L:=\operatorname{tr}A_L(\infty)+\sum_{s\geq0}\norm{\Delta L_s}^2,
\]
and similarly for \(K\).  Thus it is enough to compare the scalar square
characteristics \(Q_L\) and \(Q_K\).

First suppose \(1\leq p<2\) and put \(\alpha=p/2\).  For a martingale with
deterministic characteristics \((A,\nu)\), write
\[
   a=\operatorname{tr}A(\infty),\qquad
   U_{A,\nu}(t)
   =
   (1-e^{-ta})
   +
   \int_{\mathbb R_+\times H}(1-e^{-t\norm{x}^2})\,d\nu(s,x).
\]
The source paper's Grigelionis representation for independent-increment
martingales decomposes the jump measure into a Poisson part over the
nonatomic-in-time component of \(\nu\) and independent single jumps over the
time atoms of \(\nu\).  Consequently the Laplace transform of
\(Q=a+\sum_s\norm{\Delta L_s}^2\) has the form
\[
  \E e^{-tQ}
  =
  e^{-ta}
  e^{-U_0(t)}
  \prod_j(1-u_j(t)),
\]
where \(U_0(t)\) is the contribution of the nonatomic part and \(u_j(t)\) are
the atom-time contributions.  As in Lemma \ref{lem:laplace-square},
\[
   (1-e^{-1})\min\{1,U_{A,\nu}(t)\}
   \leq
   1-\E e^{-tQ}
   \leq
   \min\{1,U_{A,\nu}(t)\}.
\]
The assumptions imply \(U_{A_L,\nu_L}(t)\leq U_{A_K,\nu_K}(t)\) for all
\(t>0\).  The fractional-moment identity used in Lemma
\ref{lem:laplace-square} yields
\[
   \E Q_L^\alpha \lesssim_\alpha \E Q_K^\alpha .
\]

Now let \(p\geq2\) and set \(\alpha=p/2\geq1\).  For the same scalar
nonnegative independent-increment square sum, the scalar Rosenthal inequality,
applied first to finite truncations of the Poisson and atom-time
decomposition and then by monotone convergence, gives
\[
   \E Q^\alpha
   \lesssim_\alpha
   \left(a+\int\norm{x}^2\,d\nu\right)^\alpha
   +
   \int\norm{x}^{2\alpha}\,d\nu .
\]
The reverse lower bounds needed for \(K\) are elementary:
\[
   \E Q_K^\alpha\geq
   \left(\E Q_K\right)^\alpha
   =
   \left(a_K+\int\norm{x}^2\,d\nu_K\right)^\alpha
\]
by Jensen, and
\[
   \E Q_K^\alpha
   \geq
   \E\sum_s\norm{\Delta K_s}^{2\alpha}
   =
   \int\norm{x}^{2\alpha}\,d\nu_K,
\]
because \(r\mapsto r^\alpha\) is superadditive on \(\mathbb R_+\).  The
assumptions control the two upper-bound terms for \(L\) by the corresponding
terms for \(K\).  Hence \(\E Q_L^\alpha\lesssim_\alpha\E Q_K^\alpha\), and
BDG completes the proof.
\end{proof}

\section*{General Local Martingales}

\begin{proof}[Proof of Theorem \ref{thm:hilbert-all-p}]
Since Hilbert spaces are UMD, Theorem
\texttt{thm:mainforDTMgencasenananana} of \cite{Yaroslavtsev2019} gives
decoupled tangent local martingales \(\widetilde M\) and \(\widetilde N\) for
\(M\) and \(N\).  The tangent estimate
\texttt{thm:tangentgencaseUMDuhoditrotasoldat} gives
\[
  \E\sup_t\norm{M_t}^p\simequiv_p
  \E\sup_t\norm{\widetilde M_t}^p,
  \qquad
  \E\sup_t\norm{N_t}^p\simequiv_p
  \E\sup_t\norm{\widetilde N_t}^p .
\]
The same source theorem constructs the decoupled tangent process so that,
conditional on the original local characteristics, \(\widetilde M\) and
\(\widetilde N\) have independent increments and deterministic local
characteristics \(([\![M^c]\!],\nu^M)\) and
\(([\![N^c]\!],\nu^N)\), respectively.

Fix a base point \(\omega\) outside the exceptional null sets.  Characteristic
domination says
\[
   [\![N^c]\!]_\infty(\omega)\leq[\![M^c]\!]_\infty(\omega),
   \qquad
   \nu^N(\omega;\mathbb R_+\times B)
   \leq
   \nu^M(\omega;\mathbb R_+\times B)
\]
for every Borel \(B\subset H\setminus\{0\}\).  The bilinear-form order implies
\(\operatorname{tr}[\![N^c]\!]_\infty(\omega)\leq
\operatorname{tr}[\![M^c]\!]_\infty(\omega)\).  Proposition
\ref{prop:ind-incr-hilbert}, applied on this conditional probability space,
therefore gives
\[
   \E\!\left(\sup_t\norm{\widetilde N_t}^p\mid
      [\![N^c]\!],\nu^N,[\![M^c]\!],\nu^M\right)
   \leq
   C_p\,
   \E\!\left(\sup_t\norm{\widetilde M_t}^p\mid
      [\![N^c]\!],\nu^N,[\![M^c]\!],\nu^M\right).
\]
Integrating over \(\omega\) and using the tangent estimates above yields the
claimed inequality.
\end{proof}

\section*{Verification Notes}

The only genuinely new step is Lemma \ref{lem:laplace-square}.  It is
insensitive to the partition of the aggregate measure into independent
summands: all such partitions have Laplace deficits comparable to
\(\min\{1,U(t)\}\), where \(U(t)\) is determined by the aggregate law alone.
This is why rare huge jumps do not break the \(p<2\) Hilbert comparison.

The proof does not use symmetry of the original variables.  Symmetrization is
used only inside Lemma \ref{lem:hilbert-mz}, via an independent copy.  The
mean-zero assumption is exactly what makes
\(\E(Z_n-Z'_n\mid Z_1,\ldots,Z_N)=Z_n\).

The line-by-line reduction audit found one important point: characteristic
domination in \cite{Yaroslavtsev2019} controls the total spatial jump
compensator \(\nu(\mathbb R_+\times\cdot)\), not the accessible and
quasi-left-continuous compensator parts separately.  Thus the general proof
above deliberately does not claim that the accessible part of \(N\) is
dominated by the accessible part of \(M\).  Instead it uses the source's full
decoupled tangent theorem and compares the resulting independent-increment
martingales by their full characteristics in Proposition
\ref{prop:ind-incr-hilbert}.  This is exactly the source reduction, but
without imposing an unjustified componentwise domination.

A small verification script, \texttt{code/check\_laplace\_partition.py},
tests random finite partitions of the same aggregate nonnegative law and
checks the Laplace-deficit bounds used in Lemma \ref{lem:laplace-square}.

\section*{Novelty and Literature Check}

The cheap run indexes were searched for the arXiv id and for the terms
\texttt{characteristically dominated}, \texttt{aggregate domination},
\texttt{Hilbert}, \texttt{p<2}, \texttt{all p}, \texttt{Rosenthal}, and
\texttt{Laplace}.  The run already contained the broad \(c_0\) counterexample,
the Hilbert \(p=2\) packet, and the Hilbert \(p\geq2\) packet, but no Hilbert
\(1\leq p<2\) or all-\(p\) packet.  Literature/source searches identified the
source reduction in \cite{Yaroslavtsev2019}, Hilbert-valued BDG and
square-function estimates in \cite{Yaroslavtsev2018}, and Rosenthal/comparison
tools such as \cite{Pinelis2012,MontgomerySmithPruss2001}, but no exact
recorded solution of the Hilbert \(1\leq p<2\) subcase of Remark 12.10.

\section*{Limitations}

This result solves the Hilbert-valued subcase, not the full UMD-space problem.
The previous \(c_0\) counterexample shows that the unrestricted Banach-space
discrete statement in Remark 12.10 is false.  The remaining interesting
direction is therefore geometric: identify which non-Hilbert UMD spaces retain
enough square-function or tail comparability to replace Lemma
\ref{lem:hilbert-mz}.

\section*{Human Review Recommendation}

First verify Lemma \ref{lem:laplace-square}; it is short and is the decisive
new below-\(2\) mechanism.  Then verify Lemma \ref{lem:hilbert-mz}, especially
the terminal lower bound using the independent copy.  Finally verify
Proposition \ref{prop:ind-incr-hilbert}, which is the hardened replacement for
the informal accessible-jump reduction: it compares the full independent-
increment characteristics after decoupling, so no separate domination of
canonical jump parts is needed.

\begin{thebibliography}{9}

\bibitem{Yaroslavtsev2019}
I. S. Yaroslavtsev,
\emph{Local characteristics and tangency of vector-valued martingales},
arXiv:1907.11588, 2019.

\bibitem{Yaroslavtsev2018}
I. S. Yaroslavtsev,
\emph{Burkholder-Davis-Gundy inequalities in UMD Banach spaces},
arXiv:1807.05573, 2018.

\bibitem{Pinelis2012}
I. Pinelis,
\emph{Rosenthal-type inequalities for martingales in 2-smooth Banach spaces},
arXiv:1212.1912, 2012.

\bibitem{MontgomerySmithPruss2001}
S. J. Montgomery-Smith and A. R. Pruss,
\emph{Rosenthal inequalities for sums of independent random variables},
arXiv:math/9811124, 1998.

\end{thebibliography}

\end{document}
